%% Fasit til kortsett fra mattelab.no – laget automatisk av bygg.py.
%% Kompiler med lualatex.
\documentclass[11pt,a4paper]{article}
\usepackage[norsk]{babel}
\usepackage[a4paper,margin=1.2cm,bottom=1.7cm]{geometry}
\usepackage{fontspec}
\usepackage{unicode-math}
\setmainfont{TeX Gyre Pagella}
\setsansfont{TeX Gyre Heros}[Scale=MatchLowercase]
\setmathfont{TeX Gyre Pagella Math}
\usepackage{xcolor}
\definecolor{blekk}{HTML}{16213A}
\definecolor{dempet}{HTML}{566077}
\definecolor{m1}{HTML}{2F4AB0}
\usepackage{tikz}
\usepackage{enumitem}
\usepackage{qrcode}
\usepackage{fancyhdr}
\usepackage[hidelinks]{hyperref}
\pdfvariable suppressoptionalinfo \numexpr 1+2+4+8+16+32+64+128+256+512\relax
\setlength{\parindent}{0pt}
\setlength{\parskip}{0pt}
\pagestyle{fancy}
\fancyhf{}
\renewcommand{\headrulewidth}{0pt}
\fancyfoot[L]{\sffamily\footnotesize\color{dempet}mattelab.no · Potenser med samme verdi · kortsett K3 · fasit}
\fancyfoot[R]{\sffamily\footnotesize\color{dempet}Side \thepage}
\begin{document}
\begin{minipage}[t]{0.8\linewidth}\vspace{0pt}
{\sffamily\small\color{dempet}Fasit for læreren · kortsett K3 · 1T}\par\vspace{0.25em}
{\fontsize{20}{24}\selectfont\bfseries Potenser med samme verdi}\par\vspace{0.5em}
Finn fire og fire kort som har samme verdi. Regn uten kalkulator.\par\vspace{0.4em}{\small\color{dempet}Kortene er nummerert i den rekkefølgen de står i utskriften.}
\end{minipage}\hfill
\begin{minipage}[t]{0.15\linewidth}\vspace{0pt}\raggedleft
\qrcode[height=1.9cm]{https://mattelab.no/varierende/like-potenser.html}
\end{minipage}\par\vspace{0.8em}

\begin{enumerate}[leftmargin=1.6em, itemsep=0.7em]
\item \begin{minipage}[t]{\linewidth}\textsf{\small K3-04}\quad $\dfrac{1}{8}$ \\ \textsf{\small K3-05}\quad $0{,}125$ \\ \textsf{\small K3-08}\quad $2^{-3}$ \\ \textsf{\small K3-28}\quad $\left(\dfrac{1}{2}\right)^3$\end{minipage}
\item \begin{minipage}[t]{\linewidth}\textsf{\small K3-06}\quad $\dfrac{3^4}{3^4}$ \\ \textsf{\small K3-12}\quad $5^0$ \\ \textsf{\small K3-24}\quad $(-2)^0$ \\ \textsf{\small K3-32}\quad $1$\end{minipage}
\item \begin{minipage}[t]{\linewidth}\textsf{\small K3-11}\quad $10^{-2}$ \\ \textsf{\small K3-14}\quad $\dfrac{10^3}{10^5}$ \\ \textsf{\small K3-18}\quad $\dfrac{1}{100}$ \\ \textsf{\small K3-23}\quad $0{,}01$\end{minipage}
\item \begin{minipage}[t]{\linewidth}\textsf{\small K3-02}\quad $3^2$ \\ \textsf{\small K3-09}\quad $9$ \\ \textsf{\small K3-19}\quad $(-3)^2$ \\ \textsf{\small K3-29}\quad $\dfrac{3^5}{3^3}$\end{minipage}
\item \begin{minipage}[t]{\linewidth}\textsf{\small K3-10}\quad $-\dfrac{3^4}{3^2}$ \\ \textsf{\small K3-13}\quad $-3^2$ \\ \textsf{\small K3-21}\quad $-9$ \\ \textsf{\small K3-22}\quad $(-1) \cdot 3^2$\end{minipage}
\item \begin{minipage}[t]{\linewidth}\textsf{\small K3-07}\quad $32$ \\ \textsf{\small K3-25}\quad $2^2 \cdot 2^3$ \\ \textsf{\small K3-27}\quad $4^2 \cdot 2$ \\ \textsf{\small K3-31}\quad $2^5$\end{minipage}
\item \begin{minipage}[t]{\linewidth}\textsf{\small K3-01}\quad $(2^3)^2$ \\ \textsf{\small K3-15}\quad $4^3$ \\ \textsf{\small K3-26}\quad $64$ \\ \textsf{\small K3-30}\quad $2^6$\end{minipage}
\item \begin{minipage}[t]{\linewidth}\textsf{\small K3-03}\quad $0{,}5$ \\ \textsf{\small K3-16}\quad $2^{-1}$ \\ \textsf{\small K3-17}\quad $\dfrac{1}{2}$ \\ \textsf{\small K3-20}\quad $\dfrac{2^3}{2^4}$\end{minipage}
\end{enumerate}
\subsection*{Til læreren}
Fellene er lagt inn med vilje: $-3^2$ og $(-3)^2$, $2^2 \cdot 2^3$ og $(2^3)^2$, og $2^5$ og $2^6$. Be elevene skrive potensreglene de brukte, ved siden av hver gruppe.
\end{document}
